%% LyX 2.4.4 created this file.  For more info, see https://www.lyx.org/.
%% Do not edit unless you really know what you are doing.
\documentclass[english]{article}
\usepackage[T1]{fontenc}
\usepackage[utf8]{inputenc}
\setcounter{secnumdepth}{-2}
\setlength{\parindent}{0bp}
\usepackage{amsmath}
\usepackage{amsthm}
\usepackage{amssymb}

\makeatletter
%%%%%%%%%%%%%%%%%%%%%%%%%%%%%% Textclass specific LaTeX commands.
\theoremstyle{definition}
\newtheorem*{problem*}{\protect\problemname}

\@ifundefined{date}{}{\date{}}
%%%%%%%%%%%%%%%%%%%%%%%%%%%%%% User specified LaTeX commands.
\usepackage{graphicx}
\usepackage{geometry}
\newcommand{\limi}{\underset{n\rightarrow\infty}{\lim}}
\newcommand{\limxz}{\underset{x\rightarrow x_{0}}{\lim}}
\newcommand{\limfxz}{\underset{x\rightarrow x_{0}}{\lim}f(x)}
\newcommand{\limfxm}{\underset{x\rightarrow x_{0}^{-}}{\lim}f(x)}
\newcommand{\limfxp}{\underset{x\rightarrow x_{0}^{+}}{\lim}f(x)}
\newcommand{\limxm}{\underset{x\rightarrow x_{0}^{-}}{\lim}}
\newcommand{\limxp}{\underset{x\rightarrow x_{0}^{+}}{\lim}}
\newcommand{\sqrm}{M_{m\times m}(\mathbb{F})}
\newcommand{\seqa}{(a_{n})_{n=1}^{\infty}}
\newcommand{\seqx}{(x_{n})_{n=1}^{\infty}}
\newcommand{\funcdr}{f:D\rightarrow\mathbb{R}}
\newcommand{\der}{\limxz\frac{f(x)-f(x_{0})}{x-x_{0}}}
\usepackage[all]{pdfprivacy}

\makeatother

\usepackage{babel}
\providecommand{\problemname}{Problem}

\begin{document}
\title{{\Large A Problem on the Convergence of a Recursive Sequence Depending
on Its Initial Value\vspace{-2em}}}
\maketitle
\begin{problem*}
Let $\left(a_{n}\right)_{n=1}^{\infty}$ be a sequence s.t. $\forall n\in\mathbb{N}\quad a_{n+1}=a_{n}\cdot(1-a_{n})$.

Prove that $\left(a_{n}\right)_{n=1}^{\infty}$ converges if and only
if $a_{1}\in\left[0,1\right]$.
\end{problem*}
\medskip{}

\rule[0.5ex]{1\columnwidth}{1pt}

\medskip{}

\begin{proof}
We will prove both directions separately.

\medskip{}

$\left(\Rightarrow\right)$ Let us assume $\left(a_{n}\right)_{n=1}^{\infty}$
is a convergent sequence. 

1. Let $\underset{n\rightarrow\infty}{\lim}a_{n}=L$. 

Since the sequence converges, the shifted sequence $\left(a_{n+1}\right)_{n=1}^{\infty}$
must also converge to $L$.

Taking the limit on both sides of the recurrence relation and applying
sequence limits arithmetic gives
\[
L=L\cdot\left(1-L\right)\iff L=L-L^{2}\iff L^{2}=0\iff L=0
\]

2. Note the following:
\[
a_{n+1}\le a_{n}\iff a_{n}\cdot(1-a_{n})\le a_{n}\iff a_{n}-a_{n}^{2}\le a_{n}\iff-a_{n}^{2}\le0\iff a_{n}^{2}\ge0
\]
 which is true for every $n\in\mathbb{N}$ - therefore $\left(a_{n}\right)_{n=1}^{\infty}$
is a monotonically decreasing sequence.

3. Let us recall that every convergent sequence is bounded, and that
a monotonically decreasing sequence which is bounded from below converges
to the infimum of its elements.

Therefore, $0$ is the infimum of the sequence elements and in particular
$a_{1}\ge0$.

4. If $a_{1}=0$, then $a_{1}\le1$. Otherwise, $a_{1}>0$ and the
infimum provides $a_{2}\ge0$.

Substituting the recurrence relation gives $a_{1}-a_{1}^{2}\ge0\iff a_{1}\ge a_{1}^{2}\iff1\ge a_{1}$.

5. Overall, $a_{1}\in\left[0,1\right]$ as required.

\medskip{}

$\left(\Leftarrow\right)$ Assume $a_{1}\in\left[0,1\right]$ .

By step 2 of the previous direction (this argument did not use convergence),
$\left(a_{n}\right)_{n=1}^{\infty}$ is a monotonically decreasing
sequence.

Let us prove by induction that $\forall n\in\mathbb{N}\ a_{n}\in\left[0,1\right]$.

For $n=1$, the assumption provides $a_{1}\in\left[0,1\right]$.

Assume for some $n\in\mathbb{N}$ that $a_{n}\in\left[0,1\right]$.
Then,
\[
0\le a_{n}\le1\iff0\le a_{n}^{2}\le a_{n}\iff-a_{n}\le a_{n}^{2}-a_{n}\le0\iff a_{n}\ge a_{n}-a_{n}^{2}\ge0\iff a_{n}\ge a_{n+1}\ge0
\]

The induction assumption provides $1\ge a_{n}$ therefore $a_{n+1}\in\left[0,1\right]$
and the induction is complete.

Let us recall that a sequence which is monotonically decreasing and
is bounded from below is convergent, hence $\left(a_{n}\right)_{n=1}^{\infty}$
is a convergent sequence.
\end{proof}

\end{document}
